\chapter{The Grammar of Controlled Divergence}\label{ch:divergence}

The members of a variant family differ from one another in particular ways, and the ways are not all alike. Some differences cost light and nothing else. Some close off switching for the rest of the film. Some make the realizations drift apart in time. This chapter classifies the principal kinds of difference by their effects on the three distances of \Cref{ch:family} and on the switch system of \Cref{ch:switch-admissibility}. The classification is meant as a working grammar for writers: a way of knowing, before a scene is written, what a given kind of divergence will cost the family.

\section{Tonal transformation}

A tonal transformation changes how events feel without changing which events occur: comic or dramatic performance, different music, different pacing of reactions, a different emphasis in framing. The comic and dramatic realizations of one plot are the paradigm.

Image distance is usually large, because tone lives in performance and cutting, and those are on screen. Fact distance is small by design: the same things happen and the same things are revealed. The realization graph therefore stays close to complete, and tonal families are the easiest to move between, subject to the limit noted at the end of \Cref{ch:switch-admissibility}: a switch can be coherent and still be jarring. Time distance is the hidden cost. Comedy and drama keep time differently, and \Cref{ch:genre} shows that this is where tonal families are hardest to build.

\section{Expository transformation}

An expository transformation changes how much the film explains: background supplied or assumed, terms defined or used, steps shown or skipped. Beginner and expert realizations are the paradigm.

Fact distance is asymmetric. The expert realization takes its background as given from the start; the beginner realization establishes the same background as it goes. Until the beginner realization has caught up, expert viewers can move into it and beginner viewers cannot move into the expert one, whose continuation presupposes the whole background. The switch system acquires one-way transitions pointing from the realization that assumes more toward the one that explains more, and they become two-way at the point where the explanation is complete. Time distance grows with every added explanation, because explaining takes time, and \Cref{ch:depth} considers where that time can come from.

\section{Procedural transformation}

A procedural transformation shows the same action at different granularity: a repair, a surgery, a calculation, carried out in full or with familiar steps compressed. It is a special case of exposition, distinguished because what varies is not background knowledge but the resolution at which an action is depicted. Its effects on switching are the same, and it has one of its own: coarse and fine realizations of the same action can share their first and last frames exactly, which makes the action's endpoints natural rendezvous.

\section{Intensity transformation}

An intensity transformation changes sensory or emotional load: how much violence is shown, how loud, how fast, how many flashes, how severe the suspense. It is often motivated by audience sensitivity rather than taste, and it is the transformation most often proposed for age-graded versions.

When done by cutting away or framing differently, it leaves fact distance near zero: the same event occurs, less of it is seen. When done by changing what happens, it becomes a causal transformation and inherits that kind's costs. \Cref{ch:unfiltered} applies here with force. A lower-intensity realization does not protect anyone from the higher-intensity one on the same screen, because the composite contains both.

\section{Perspective transformation}

A perspective transformation follows different characters, or places the camera differently, through the same events. Its effect on fact distance depends on a single question: does one perspective show something the other does not? Two views of a conversation from opposite sides differ only in salience. A view that follows a character into a room the other view never enters differs in what it establishes. Perspective transformations can therefore be as cheap as tonal ones or as consequential as causal ones, and the difference is decided by information, not by camera position. \Cref{ch:viewpoint} develops this.

\section{Causal transformation}

A causal transformation changes what happens: a character decides differently, an event occurs in one realization and not the other. It is the most powerful kind of divergence and the most expensive.

Fact distance grows with every consequence. Each realization establishes literals the other contradicts, and the realization graph thins accordingly. Unless the realizations reconverge, bringing their established facts back into agreement on everything later continuations depend on, switching is closed for good. Reconvergence is possible but has a narrative price: the story must supply reasons why different decisions led to the same place. Stories built this way tend toward fatalism, in which choices change the route and not the destination, and a writer should recognize that the structure imposes the theme.

\section{Terminal transformation}

A terminal transformation is a causal transformation placed at the end: realizations that share everything until a final branch and then resolve differently. It avoids the price of reconvergence by never paying it. The realization graph is complete until the branch and empty after it, and image distance is zero for most of the running time, which makes terminal families nearly free in light.

The cheapness is partly deceptive. An ending changes the meaning of what came before, so a family with different endings is not really identical until the branch. The same scenes are seen differently depending on where they lead, and if the endings are to feel earned, the shared portion must contain the seeds of each. \Cref{ch:endings} distinguishes the ways of doing this.

\section{Localization transformation}

A localization transformation changes language, performers, cultural references, signage, or setting for different audiences. When it changes only language, it is close to tonal: same events, different surface. When it changes jokes, customs, or performers, it can alter what is established about characters and their world, and it acquires fact distance of its own. It also raises questions of authorship that the other transformations do not, because it asks one culture's realization to stand in for another's. \Cref{ch:local} takes these up.

\section{Composition and interference}

Transformations are rarely used alone, and they interfere with each other. A family that varies both tone and explanation does not get four independent realizations for free. Comedy compresses and drama lingers, so the time available for explanation differs between the tonal variants, and an expert comic realization may leave no room for what a beginner dramatic realization needs. An ending that suits a comic realization may be unearned in a dramatic one. The nominal count of a family, the product of its dimensions, overstates what it actually contains whenever the dimensions interact.

The three distances are analytically independent, in the sense that each can be varied while the others are held fixed. They are not operationally independent, in the sense that real transformations rarely touch only one. Comedy alters timing as well as image. Explanation introduces facts as well as time. An alternate ending can demand different earlier imagery to seed it. If $\Delta d_i(T)$ is the change in distance $i$ produced by a transformation $T$, then in general
\begin{equation}
\Delta d_i(T_a\circ T_b)\ \neq\ \Delta d_i(T_a)+\Delta d_i(T_b),
\end{equation}
and variants that are cheap when designed separately can become expensive when combined.

Even so, the three distances give a way to anticipate the interference to first order. Tonal and localization changes mostly spend image distance. Expository and procedural changes mostly spend time distance, with one-way fact distance. Causal and terminal changes spend fact distance. Two transformations that spend the same distance compete for it; two that spend different distances can often be combined. A comic and dramatic family with beginner and expert realizations competes for time, and must be planned on the clock. A comic and dramatic family with two endings spends image distance throughout and fact distance only at the end, and combines easily, unless the endings need different seeds in the comic and dramatic realizations. The eight-way prototype of Appendix~E, crossing genre, depth, and ending, exists precisely to measure these couplings.

\section{A writer's summary}

The grammar reduces to three questions a writer can ask of any proposed divergence. Will it show different pictures? Then it costs light. Will it establish different facts? Then it costs switching, until the facts are brought back into agreement. Will it take different amounts of time? Then it costs synchronization, until the time is repaid. Every difference in a variant family pays in at least one of these currencies, and the art of designing a family is largely the art of choosing which.
